\documentclass[11pt,a4paper]{article}
\usepackage[utf8]{inputenc}
\usepackage[T1]{fontenc}
\usepackage{amsmath,amssymb}
\usepackage{booktabs}
\usepackage{geometry}
\geometry{left=2.2cm,right=2.2cm,top=2.2cm,bottom=2.2cm}
\usepackage{url}
\usepackage{hyperref}
\hypersetup{colorlinks=true,linkcolor=blue,citecolor=blue,urlcolor=blue}
\usepackage{enumitem}

\title{\textbf{BRICK: Branch-Rewired Inference with Composable Knowledge}\\
\large A Negative-Result Study of 2D Layer Topologies for Dynamic Activation\\
\large in Frozen Small Language Models}
\author{Weikai Yue, Hy 4 flash}
\date{Independent research report --- Draft v0.1 --- 2026-11-25}

\begin{document}
\maketitle

\begin{center}
\textit{Follow-up to Amoeba \quad $\cdot$ \quad CPU-only empirical study \quad $\cdot$ \quad All results, including negative, reported}
\end{center}

\begin{abstract}
We ask whether the one-dimensional layer chain of a pretrained transformer can be
rewired into a two-dimensional topology --- parallel heterogeneous branches fused
before a shared tail --- to enable dynamic activation without retraining the backbone.
Motivated by the FinFET transition (changing geometry, not material) and by our earlier
finding that middle layers of a 1B model have near-zero ablation cost, we build BRICK:
a frozen MiniCPM5-1B whose 24 layers are rearranged as an 8-layer prelude, two parallel
branches of 8 and 6 layers, and a 2-layer coda, with only a lightweight fusion module
and router trained. Across a complete CPU-only experimental chain we find:
(1) the 2D rearrangement itself is \emph{not} capability-destroying --- a trained
fusion recovers the ``both-branches'' cross-entropy from 4.22 to 4.04;
(2) however, \emph{skipping} a branch carries an irreducible cost of $+0.40$--$0.90$ nats,
because the skipped branch's information is genuinely absent and cannot be reconstructed
by fusion; (3) exposing the fusion to single-branch inputs during training makes results
\emph{worse}, not better; (4) a Gumbel-Softmax router that balances branch usage on
synthetic data collapses back to the deep branch on the real model. We report the full
negative-result chain, the layer/block importance ablation that motivated the topology,
and the boundary conditions under which 2D rewiring is and is not viable. We claim no
novelty for the mechanisms (parallel branches, fusion, routers, early exit); the
contribution is the empirical boundary map and the honest negative result.
\end{abstract}

\section{Introduction}

Amoeba established that dynamic computation in frozen small models is bounded by
architecture and harness: speculative decoding works on full-attention dense models,
while early exit fails on pretrained backbones. This study asks a different question:
what if we change the \emph{geometry} of the layer stack itself?

The analogy is FinFET: the transistor did not change material, it changed shape
(planar channel $\to$ 3D fin) to regain control and density. For transformers, the
universal assumption is that layers form a 1D chain $0,1,\ldots,L{-}1$. We test a 2D
execution graph: several layers are executed in parallel as branches from a shared
hidden state and then fused. This is ``composable knowledge'': the pretrained weights
are kept frozen, and only the composition is learned.

We stress from the outset that parallel branches, mixture-of-depths, fusion and
routers are all published mechanisms (Multi-Branch Transformer, Mixture-of-Depths,
GateSkip). Our study is not a new
mechanism; it is a controlled empirical test of whether these mechanisms, applied to
a \emph{frozen} small model with a \emph{heterogeneous} 2D topology, actually deliver
the promised dynamic activation. The answer, after a complete experimental chain,
is largely negative --- and that is the contribution.

\section{Background and Motivation}

\subsection{Earlier Amoeba findings}
\begin{enumerate}[leftmargin=*]
  \item Token-level early exit on pretrained MiniCPM5-2B failed in 16 experiments;
        the last layers alone contribute $\sim$1.3 nats CE and there is no redundant
        depth to skip; structural ceiling $\sim$1.16$\times$.
  \item Speculative decoding (Amoeba-DA) gives 2.015$\times$ on 2B but depends on a
        verification engine incompatible with GQA/sliding-window KV caches; all Qwen
        speedups were withdrawn.
  \item Layer-importance analysis on 1B showed a U-shaped profile: layers 0 and 4 are
        ``cornerstone'' layers ($\Delta$CE $+6.5$, $+5.6$), layers 22--23 are also
        critical, while layers 5--21 have near-zero ablation cost.
\end{enumerate}
The third finding is the motivation for BRICK: if middle layers are redundant, can we
arrange them as parallel branches and activate only some of them?

\subsection{Why LayerSkip works and our situation does not}
LayerSkip (Elhoushi et al., ACL 2024) also exploits the fact that early/middle layers can be
skipped --- but it \emph{trains for it}: layer dropout (low rate early, high rate late)
plus an early-exit loss on a shared LM head make intermediate representations directly
decodable. BitSkip (Bhuvaneswaran \& Liu, 2026) independently reproduces this: a shared LM head
only becomes a useful intermediate exit after from-scratch training (per-layer exit PPL
drops from 12,418 at layer 10 to 216 at layer 11; the final-layer cliff in both studies).
We have a pretrained backbone and no early-exit loss: redundancy exists under
\emph{ablation} (later layers compensate) but not under \emph{readout} (intermediate
hidden states cannot be decoded). BRICK's learned fusion is an attempt to bridge this
gap with a lightweight module instead of full retraining.

\section{Method}

\subsection{2D topology}
For MiniCPM5-1B (24 layers, hidden 1536), we use:
\[
  \underbrace{L_{0..7}}_{\text{prelude}(8)}
  \;\longrightarrow\;
  \Big[\underbrace{L_{8..15}}_{\text{branch}_0(8)},
       \underbrace{L_{16..21}}_{\text{branch}_1(6)}\Big]
  \;\longrightarrow\;
  \underbrace{\text{Fusion}}_{\text{learned}}
  \;\longrightarrow\;
  \underbrace{L_{22..23}}_{\text{coda}(2)}
  \;\longrightarrow\;
  \text{LM head}.
\]
The two branches consume the same prelude output and execute in parallel; each uses
\emph{distinct original layer indices}, so in HuggingFace's dynamic KV cache the 24
cache slots are each written exactly once per token and no cache conflict arises.
The coda always contains the final layer, since our earlier work shows it is the
readout interface.

\subsection{Fusion}
The default fusion is a residual correction anchored at the branch mean:
\[
  h_{\text{out}} = \frac{1}{K}\sum_k b_k \;+\; \alpha\, W\!\left(\mathrm{concat}(b_1,\ldots,b_K)\right),
\]
with $W$ zero-initialized and $\alpha=1$ --- the start point is exactly the branch
mean (no representation damage), while $W$ retains full gradient. Two alternatives
were tested and rejected (see \S5): prelude-anchored residual
($h_{\text{pre}}+\alpha W(\cdot)$) and an unanchored ``direct'' projection.

\subsection{Router}
A Gumbel-Softmax router (temperature annealed $1.0\to0.1$, straight-through)
selects which branch(es) to execute per token/segment. Three oracle definitions were
tested: (a) representation distance to a 1D teacher, (b) per-branch CE, (c)
\emph{ablation aware}: the per-token cost of skipping each branch, i.e.\ the genuinely
early-exit-compatible signal. Only fusion and router parameters are trained
($<$10M); the 2GB backbone is frozen.

\section{Experimental Setup}

All experiments run on a CPU-only machine (Intel i7-1260P, 15.6GB RAM, PyTorch 2.15
CPU, transformers 5.24.1). The 1B model loads in bf16 (RSS 2.36GB); one forward pass
takes $\sim$0.3\,s, one training step (14 branch layers + fusion + coda + 130k-vocab
LM head) takes 1.6--15\,s depending on whether CE is computed. We therefore use small
settings (3--6 windows, seq 16, 12--50 steps) to identify \emph{trends and boundaries},
not to claim final accuracy. Critically, two bugs that would have destroyed a remote
GPU run were caught by these local tests: transformers 5.x requires
\texttt{position\_embeddings=(cos,sin)} rather than \texttt{position\_ids}, and the
fusion parameters must match the backbone dtype (bf16 vs fp32).

\section{Results}

\subsection{The 2D rearrangement is not itself destructive}
With untrained fusion, the both-branch CE is 4.219. After 50 fusion steps at
LR $3\times10^{-5}$, it drops to \textbf{4.041} (MSE $1.42\to0.03$). The gap between
2D and 1D is learnable by a lightweight module --- the ``composable knowledge''
premise holds for the \emph{fused} path.

\subsection{Skipping a branch has an irreducible cost}
In the full chain (trained fusion + ablation oracle + Gumbel router), the router's
end-to-end CE is \textbf{4.4435}, i.e.\ $+0.40$ above both (4.041). The ablation
oracle gives skip costs: both 4.04, skip branch$_0$ 4.44, skip branch$_1$ 4.83
--- skipping either branch costs 0.4--0.8 nats. This is the central negative result:
the skipped branch's information is genuinely absent; fusion cannot reconstruct it.

\subsection{Exposing fusion to single-branch inputs makes it worse}
We trained fusion under three input strategies for 12 steps: both-only, mixed
(both/skip$_0$/skip$_1$), single-only. End-to-end router CE:
both-only 5.16, mixed 5.37, single-only 5.41; both CE itself degraded from 4.267 to
4.420 under single-branch training. Forcing fusion to ``guess'' missing information
both fails to guess and harms normal fusion. This closes the most promising remedy.

\subsection{Router behavior}
On synthetic tiny models, a plain softmax router collapses to a single branch
($[1.0,0.0]$); Gumbel-Softmax balances it to $[0.48,0.52]$ without any extra entropy
or load-balancing term. On the real 1B model, however, the Gumbel router still ends at
$[0.09,0.91]$ (all three oracles converge to the same distribution and the same CE),
because the real branch quality gap dominates the random exploration. The oracle
definition matters little once the backbone is frozen and fused once.

\subsection{Layer and block importance}
\begin{table}[h]
\centering
\begin{tabular}{lcc}
\toprule
Ablation (1D sequential) & $\Delta$CE & \\ \midrule
layer 0 & $+6.51$ & cornerstone \\
layer 4 & $+5.62$ & cornerstone \\
layer 22 / 23 & $+5.02$ / $+3.62$ & readout \\
layers 5--21 & $\approx 0$ & redundant \\
\bottomrule
\end{tabular}
\caption{1D layer ablation: U-shaped importance (real 1B, 6 windows).}
\end{table}
In the BRICK 2D topology, skipping prelude layers 4--7 is catastrophic
(CE $3.79\to15.66$, $\Delta=+11.87$), while skipping a branch costs only
$+0.31$--$0.36$. The topology therefore places the critical layers in the
prelude/coda and the redundant layers in the branches --- a design later justified
by ablation. Dynamic activation is legally confined to the branch pair.

\subsection{Scale dependence}
Amoeba's 2B model has the last layer dominating ($\sim$1.3 nats); in BRICK's 1B the
first layer dominates ($+6.51$) and the final layer contributes $+3.62$. Layer
importance is scale-dependent, so a 2D topology optimal at one scale may not transfer.

\section{Negative Results and Boundaries}
\begin{enumerate}[leftmargin=*]
  \item \textbf{Branch skipping is not free in 2D}: $1D$ ablation suggested $+0.07$;
        in the 2D fused topology the real cost is $+0.4$--$0.9$ nats.
  \item \textbf{Fusion cannot reconstruct missing information} (\S5.3).
  \item \textbf{Gumbel balancing is scale-dependent}: works on tiny random models,
        fails on real 1B.
  \item \textbf{Offline oracle mismatch}: an ablation oracle computed with untrained
        fusion locks all routers to one branch; oracle quality is only meaningful after
        fusion training.
  \item \textbf{CPU limits}: full convergence (hundreds of steps) and downstream task
        evaluation require a GPU; the reported numbers are trends/boundaries, not
        final accuracy.
  \item \textbf{A6 (depth interpolation)} would preserve capability by evaluating
        both branches but therefore abandons the acceleration goal; it is left as an
        explicitly non-accelerating option.
\end{enumerate}

\section{Related Work (positioning, not novelty claims)}
\begin{itemize}[leftmargin=*]
  \item \textbf{Layer importance / pruning}: Investigating Layer Importance in LLMs
        (BlackboxNLP 2024), The Curse of Depth in LLMs (NeurIPS 2025),
        Reassessing Layer Pruning (ICLR 2026), No Free Swap (2026). Our 1D ablation
        reproduces their U-shaped / cornerstone findings; not novel.
  \item \textbf{Early exit}: LayerSkip (ACL 2024), BitSkip (2026). Both require
        from-scratch training with an early-exit loss; we test a frozen backbone.
  \item \textbf{2D / parallel topologies}: Multi-Branch Transformer (2026) and
        DAGFormer (2026) use branches/DAG connections but train from scratch;
        mHC / TeamFormer change training. Our difference is the frozen backbone and
        the negative-result boundary, not the mechanism.
  \item \textbf{Speculative decoding}: Amoeba-DA (prior work), Medusa, EAGLE, JetSpec.
\end{itemize}

\section{Conclusion}
BRICK asked whether a pretrained small model's 1D layer chain can be rewired into a
2D parallel-branch topology for dynamic activation. The answer is nuanced and mostly
negative: the fused path is fully recoverable by a lightweight module (CE 4.22$\to$4.04),
but \emph{activating only one branch} is not --- the missing branch carries real
information whose loss costs 0.4--0.9 nats and cannot be trained away. We release the
complete code, all experiments, and the full negative-result chain on GitHub. The
study's value lies in the boundary map and in demonstrating, once again, that a plausible
mechanism (parallel branches + router) does not automatically become an accelerator
when applied to a frozen model.

\end{document}
